\chapter{Ectopic Pregnancy}
\index{Ectopic Pregnancy}
\label{sec:Ectopic Pregnancy}

\section{Overview}
Ectopic pregnancy occurs when a fertilized ovum implants outside the uterine cavity, most commonly in the fallopian tube (95\%), but also in the cervix, ovary, abdominal cavity, or previous cesarean scar. The incidence is 1--2\% of all pregnancies. This pregnancy-related condition requires careful monitoring to ensure the safety of both mother and baby. Early detection and management are essential. This condition affects a significant number of people worldwide and can range from mild to severe in presentation. Early recognition and appropriate management are essential for optimal outcomes. A thorough understanding of the underlying mechanisms guides treatment decisions. Patient education and self-management play important roles in long-term care.
\section{Causes}
This condition happens because of impaired tubal function or anatomy preventing normal embryo transport. The underlying mechanisms involve complex interactions between genetic predisposition and environmental factors that disrupt normal physiological processes. The underlying mechanisms involve complex interactions between genetic predisposition and environmental factors. Disruption of normal physiological processes leads to the characteristic features of the condition. Multiple contributing factors may act synergistically to trigger or worsen the condition. Understanding the specific cause helps guide targeted treatment approaches.
\section{Risk Factors}
Genetic predisposition and family history Advancing age Lifestyle factors (diet, exercise, smoking, alcohol) Environmental exposures Underlying medical conditions Medication use Stress and psychological factors Socioeconomic factors

\begin{itemize}[noitemsep]
  \item Genetic predisposition and family history
  \item Advancing age
  \item Lifestyle factors (diet, exercise, smoking, alcohol)
  \item Environmental exposures
  \item Underlying medical conditions
  \item Medication use
\end{itemize}
\section{Signs and Symptoms}
\begin{itemize}[noitemsep]
  \item You may experience amenorrhea, vaginal bleeding, and unilateral lower abdominal pain
  \item Ruptured ectopic pregnancy presents with sudden severe pain, syncope, involving bleeding shock, and peritonitis.
  \item Pain or discomfort in the affected area
  \item Functional impairment affecting daily activities
  \item Changes in appearance or function of affected tissues
  \item Systemic symptoms may include fatigue, fever, or weight changes
  \item Symptoms may develop gradually or appear suddenly
\end{itemize}
\section{Progression and Stages}
The condition may progress through different stages of severity over time. Early stages often involve mild or subtle symptoms that may be overlooked. Moderate stages involve more noticeable symptoms affecting daily function. Advanced stages may involve significant impairment and complications.
\section{Complications}
\begin{itemize}[noitemsep]
  \item Disease progression leading to worsening symptoms
  \item Secondary infections or injuries
  \item Side effects of treatment
  \item Reduced quality of life
  \item Emotional and psychological impact
\end{itemize}
\section{Diagnosis}
\begin{itemize}[noitemsep]
  \item Doctors diagnose this using serum $\beta$-hCG and transvaginal ultrasound (empty uterine cavity, adnexal mass, free fluid)
  \item A discriminatory zone of $\beta$-hCG >2000 IU/L with no intrauterine pregnancy suggests ectopic pregnancy.
  \item Comprehensive medical history and physical examination
  \item Laboratory tests to assess organ function
  \item Imaging studies to visualize affected structures
  \item Specialized testing depending on the affected system
  \item Consultation with relevant specialists
\end{itemize}
\section{Prevention}
\begin{itemize}[noitemsep]
  \item Maintain a healthy lifestyle with balanced diet and exercise
  \item Avoid known risk factors
  \item Attend regular medical check-ups for early detection
  \item Follow recommended screening guidelines
\end{itemize}
\section{Treatment Options}
\begin{itemize}[noitemsep]
  \item The right treatment depends on stability: stable patients with small ectopic may be treated with methotrexate
  \item Unstable patients require using small incisions and a camera salpingectomy or salpingotomy.
  \item Medications to manage symptoms and address underlying mechanisms
  \item Lifestyle modifications and self-care measures
  \item Physical therapy and rehabilitation
  \item Surgical intervention when indicated
  \item Regular monitoring and follow-up care
\end{itemize}
\section{Living with It}
\begin{itemize}[noitemsep]
  \item Follow your treatment plan as prescribed
  \item Maintain regular follow-up with your healthcare provider
  \item Make necessary lifestyle adjustments
  \item Seek support from family, friends, or support groups
  \item Stay informed about your condition
\end{itemize}
\section{Outlook}
\begin{itemize}[noitemsep]
  \item The outlook is generally positive with early diagnosis
  \item Ruptured ectopic pregnancy is life-threatening.
\end{itemize}
